% GENERATED FILE. Edit latin/pronunciation-reference.md, then run npm run generate:books.

\chapter*{How Latin sounds --- a reference}
\addcontentsline{toc}{chapter}{How Latin sounds (reference)}
\markboth{Pronunciation}{Pronunciation}

A \textbf{reference}, not a chapter — the lessons teach the sounds through real words. This page gathers the classical scheme (the one this course uses).

\section*{The three facts}

\begin{enumerate}
\raggedright
\setlength{\itemsep}{0.35em}
  \item \textbf{Every letter is sounded, always the same way.} No silent letters and, in the classical scheme, no \textquotedblleft{}soft\textquotedblright{} consonants: \emph{c} and \emph{g} are \textbf{always hard} (\emph{Cicero} = \emph{KIH-keh-ro}, not \textquotedblleft{}SIS-er-o\textquotedblright{}), \emph{v} is \textbf{w}, \emph{s} is always hissed (never \textquotedblleft{}z\textquotedblright{}), \emph{r} is trilled.
  \item \textbf{Vowels come long or short, marked by a macron.} ā ē ī ō ū are \emph{held longer} than a e i o u — same quality, more time. Length can change meaning (\emph{malum} \textquotedblleft{}evil\textquotedblright{} vs. \emph{mālum} \textquotedblleft{}apple\textquotedblright{}), which is why teaching texts print macrons Romans never wrote.
  \item \textbf{Diphthongs:} \emph{ae} = \textquotedblleft{}eye\textquotedblright{} (\emph{Caesar} = \emph{KYE-sar}), \emph{au} = \textquotedblleft{}ow,\textquotedblright{} \emph{oe} = \textquotedblleft{}oy\textquotedblright{}; otherwise adjacent vowels each keep their value.
\end{enumerate}

\section*{Worth flagging}

\begin{itemize}
\raggedright
\setlength{\itemsep}{0.35em}
  \item \emph{\textbf{v\emph{ = w, }i\emph{ = y\textbf{ as consonants: \emph{vīnum} (\textquotedblleft{}wine\textquotedblright{}) = \emph{WEE-noom}, \emph{Iulius} = \emph{YOO-lee-oos}. The \textquotedblleft{}v as vee / soft c\textquotedblright{} you may have heard is \emph{later} Church Latin.}}}}
  \item \textbf{Stress is predictable:} two-syllable words stress the first; longer words stress the second-to-last syllable if it is long, else the one before. No accent marks are needed once you hear vowel length.
  \item \textbf{No \emph{j}, \emph{u/v} split, or \emph{w}:} classical Latin wrote \textbf{I} and \textbf{V} for both vowel and glide; \emph{j}, \emph{u}, \emph{w} are medieval/modern additions.
\end{itemize}

\section*{Latin in the family}

Latin is \textbf{Italic}, a branch of \textbf{Indo-European} — cousin to Greek, to the Germanic languages (English among them), and, further off, to Sanskrit. It was the speech of Rome, then of an empire, then of the medieval church and the learned world; as it fragmented in daily use it became Spanish, Portuguese, French, Italian, Romanian, and the rest. This curriculum treats Latin as a \textbf{taproot}: nearly every word has both English and Romance descendants, traced in the lessons — and its sister taproot, Sanskrit, shares the same Indo-European negative (\emph{nōn} / \emph{na}) and the family of \emph{mother/father} words.
